\section{Discussion}
\label{sec:discussion}

\paragraph{What the results support.}
Three findings hold across the methods used here. First, \sys{} writes decisive reports. The rater found a firm position far more often than in either baseline, and the lexical measurements agree, with fewer hedges on nine of ten motions and more commitment markers on all ten. Second, \sys{} writes compact reports organized around a claim and cites densely within them. Third, its recursive composition, as implemented in the evaluated version, did not preserve the link between each claim and the source that supports it. The rater's higher scores for supported claims and evidence use reflect the first two findings, namely consistent citation in a compact report, rather than correct attribution.

\paragraph{The stance result depends partly on the prompt.}
The report prompt instructs the model to argue for the root answer and not to hedge, so a high rate of firm positions would be expected from the prompt alone. The decomposition contributes a root answer in a definite format for the prompt to argue from, but the experiments do not separate these two contributions. The analysis in \Cref{sec:analysis} adds a further observation. Titles and executive summaries commit to a position, but half of the conclusions qualify or condition it. A likely reason is that the final writing step draws on a graph whose branches cover both sides of the question, so the balanced material reappears at the end. Three ablations would separate these effects. The first keeps the graph and removes the stance instruction, the second keeps the instruction and removes the graph, and the third requires the conclusion to agree with the title.

\paragraph{Why provenance fails in recursive composition.}
A citation number is a pointer into a list of sources. The evaluated pipeline rebuilt that list at every level of the graph but copied the pointers without changing them, so each pointer came to refer to a different source. The same risk arises in any system that composes cited text in stages, for example by summarizing summaries or by merging the outputs of several agents. The error is invisible both to a reader and to a rubric that counts citation markers, because the report still looks well cited. Two design rules prevent it. Each source should carry one identifier for the whole run, and no model should refer to an intermediate result by an identifier that the final text cannot resolve to a source. A lexical audit such as the one in \Cref{sec:citations} offers an inexpensive check that both rules hold, and attribution measures based on entailment models \citep{rashkin2023ais, gao2023alce} offer a stronger one.

\paragraph{Scale and raters.}
Ten motions and six transcripts are too few for precise estimates. The human scores come from one rater who is also an author and who, although blind to the system, was not blind to the study. The debates were judged by a small model applying a rubric from an unpublished course project, and two rubric measures had to be dropped because their records were unusable. The paired measurements in \Cref{sec:analysis} do not depend on a rater, but they are lexical proxies, and the stance labels in \Cref{tab:stance} come from a single reading.

\paragraph{Uncontrolled factors.}
The three systems differ in their models, search backends, and prompts, so the comparison measures complete systems rather than the effect of decomposition alone. STORM was built to write balanced encyclopedia articles and is evaluated here on a task it was not designed for, so its low count of firm positions reflects its design rather than a failure. Deep Research's reports were produced by a commercial product whose internals and model version cannot be controlled. A controlled study would hold the language model and the search API fixed across a flat search agent, STORM, and \sys{}, would include the ablations described above, and would measure attribution with the test in \Cref{sec:citations} alongside human judgment.

\paragraph{Cost.}
A \sys{} run makes on the order of four model calls per leaf and a few per parent (\Cref{sec:method}). Its cost therefore grows linearly with the size of the graph and is several times that of a report written in a single pass. Much of what it retrieves is never used, since only a median of \atUsedShare\% of the listed sources are cited in the body of the report. Removing sources that no claim uses before the final step would shorten the bibliography without changing the argument.
